\documentclass[11pt]{article}
\usepackage[a4paper,margin=1in]{geometry}
\usepackage{fontspec}
\usepackage[boldfont=false]{xeCJK}
\setCJKmainfont{FandolSong-Regular.otf}
\usepackage{amsmath}
\usepackage{amssymb}
\usepackage{array}
\usepackage{booktabs}
\usepackage{graphicx}
\usepackage{adjustbox}
\usepackage{enumitem}
\setlist[itemize]{topsep=2pt,partopsep=0pt,parsep=0pt,itemsep=2pt,leftmargin=1.4em}
\setlist[enumerate]{topsep=2pt,partopsep=0pt,parsep=0pt,itemsep=2pt,leftmargin=1.6em}
\usepackage{parskip}
\usepackage{fvextra}
\fvset{breaklines=true,breakanywhere=true,fontsize=\small}
\setlength{\parindent}{0pt}

\begin{document}

\begin{center}
  {\LARGE\bfseries Biotechnology and genetic modification}\\[0.35em]
  {\large Igcse Biology}
\end{center}
\vspace{0.6em}

\section*{Why bacteria are useful}

\textbf{Biotechnology} 生物技术 uses living things (or their \textbf{enzymes} 酶) to make useful products. \textbf{Bacteria} 细菌 are especially useful because:

\begin{itemize}
\item they reproduce very fast (a high \textbf{rate} 速率 of reproduction).

\item they can make complex \textbf{molecules} 分子.

\item \textbf{(Supplement)} there are few \textbf{ethical} 伦理 concerns about growing them, and they contain \textbf{plasmids} 质粒 (small DNA rings that are easy to work with).

\end{itemize}

\section*{Biotechnology in action}

\subsection*{Yeast: biofuels and bread}

\textbf{Yeast} 酵母 carries out \textbf{anaerobic respiration} 无氧呼吸 (without oxygen), making \textbf{ethanol} 乙醇 and carbon dioxide.

\begin{itemize}
\item the ethanol can be used as a \textbf{biofuel} 生物燃料.

\item in bread-making, the carbon dioxide makes the dough rise.

\end{itemize}

\subsection*{Enzymes in industry}

\begin{itemize}
\item \textbf{pectinase} 果胶酶 breaks down cell walls to release more juice from fruit, giving more and clearer juice.

\item biological \textbf{washing powders} 洗衣粉 contain enzymes (such as proteases and lipases) that digest stains, even at lower temperatures.

\item \textbf{(Supplement)} \textbf{lactase} 乳糖酶 breaks down \textbf{lactose} 乳糖 to make lactose-free milk, for people who cannot digest lactose.

\end{itemize}

\subsection*{Fermenters (Supplement)}

A \textbf{fermenter} 发酵罐 is a large tank used to grow bacteria or fungi to make useful products, such as \textbf{insulin} 胰岛素, \textbf{penicillin} 青霉素 and \textbf{mycoprotein} 真菌蛋白. The conditions inside must be carefully controlled: \textbf{temperature} 温度, pH, \textbf{oxygen} 氧气, the supply of nutrients, and the removal of waste products.

\section*{Genetic modification}

\textbf{Genetic modification} 基因改造 means changing an organism's \textbf{genetic material} 遗传物质 by removing, changing or inserting individual \textbf{genes} 基因.

\subsection*{Making a human protein in bacteria (Supplement)}

For example, to make a human \textbf{protein} 蛋白质 (such as insulin) in bacteria:

\begin{enumerate}
\item cut the human gene out of human DNA using \textbf{restriction enzymes} 限制酶, which leave \textbf{sticky ends} 黏性末端.

\item cut open a bacterial plasmid with the \textbf{same} restriction enzymes, giving matching sticky ends.

\item join the human gene into the plasmid using DNA \textbf{ligase} 连接酶, making a \textbf{recombinant plasmid} 重组质粒.

\item put the recombinant plasmid into a bacterium.

\item the bacteria multiply and \textbf{express} 表达 the human gene, making the human protein.

\end{enumerate}

\subsection*{Examples of genetic modification}

Genetic modification is used to:

\begin{itemize}
\item put human genes into bacteria to make human proteins (such as insulin).

\item put genes into \textbf{crop plants} 农作物 to give \textbf{resistance} 抗性 to \textbf{herbicides} 除草剂.

\item put genes into crops to give resistance to insect pests.

\item put genes into crops to improve their food value.

\end{itemize}

\textbf{(Supplement)} GM crops such as soya, maize and rice can give higher yields and better nutrition, but some people worry about effects on health, on wild species, and about the cost of GM seeds.

\section*{Exam tips}

\begin{itemize}
\item Bacteria are used because they reproduce fast, make complex molecules, and contain plasmids.

\item Yeast (anaerobic respiration) → ethanol (biofuel) and carbon dioxide (bread). Enzymes: pectinase (fruit juice), proteases and lipases (washing powders), lactase (lactose-free milk).

\item Fermenters make insulin, penicillin and mycoprotein; control temperature, pH, oxygen, nutrients and waste.

\item Genetic modification: cut a gene with restriction enzymes (sticky ends) → join it into a plasmid with ligase (a recombinant plasmid) → put into bacteria → express the gene to make the protein.

\item GM crops can resist herbicides or insect pests, or have better nutrition.

\end{itemize}


\end{document}
